\title{Direct Preference Optimization: Your Language Model is Secretly a Reward Model}

\begin{abstract}
Although large-scale unsupervised language models (LMs) capture extensive general knowledge and develop some capacity for reasoning, exerting fine-grained control over their outputs is challenging, primarily because their training lacks explicit supervision. Current approaches aiming to improve the steerability of LMs generally involve collecting human judgments comparing pairs of model outputs, followed by fine-tuning the base model to better match these preferences—most often through reinforcement learning from human feedback (RLHF). Nonetheless, RLHF remains a complex and sometimes unstable method: it involves training a reward model to emulate human preferences, then employing reinforcement learning to optimize the LM towards this reward while attempting to maintain fidelity to the original model parameters. In this work, we propose an alternative parameterization for the reward model within the RLHF framework that permits direct extraction of the corresponding optimal policy in closed form. This advancement enables us to address the conventional RLHF setup using a straightforward classification loss alone. We introduce the resulting technique as \textit{Direct Preference Optimization} (DPO), which offers stability, strong performance, and reduced computational burden, as it obviates the need to sample from the language model during fine-tuning or to engage in extensive hyperparameter adjustment. Experimental results demonstrate that DPO is capable of adapting LMs to align with human preferences on par with or better than prevailing approaches. Remarkably, DPO outperforms PPO-based RLHF in sentiment control tasks, and achieves equal or superior results in summarization and single-turn dialogue quality, all while providing a far simpler training and implementation process.
\end{abstract}

\section{Introduction}
Unsupervised large language models (LMs) trained on massive text corpora have demonstrated remarkable emergent abilities~\citep{chowdhery2022palm, brown2020language, touvron2023llama, bubeck2023sparks}. Nonetheless, these models learn from human-generated data encompassing a vast array of intentions, standards, and expertise. As a consequence, not every goal or competency present in the data is necessarily one we wish to replicate. For instance, while it is beneficial for an AI-based programming assistant to \textit{recognize} frequent coding errors in order to address them, we prefer the model to favor the rare yet high-quality code examples within its training corpus during code generation tasks. Likewise, although it is valuable for the model to \textit{know about} a misconception held by half of the population, it is undesirable for the model to reproduce that misconception as truth in half of its relevant outputs. In sum, it is imperative to distinguish and select the model’s \emph{intended behaviors and outputs} from its extensive \textit{repertoire of knowledge and skills} to ensure that AI systems are effective, safe, and align with user values \citep{ouyang2022training}. Whereas prevailing techniques utilize reinforcement learning (RL) to shape LMs in accordance with human judgments, we demonstrate that the RL-based objective frequently adopted in such settings can be achieved precisely through a straightforward binary cross-entropy formulation, thereby streamlining the preference optimization process.

In broad terms, current techniques embed preferred behaviors into language models by utilizing carefully selected collections of human preference data, reflecting the actions humans deem safe and useful. This process of learning from preferences is performed following an initial phase of large-scale, unsupervised pre-training on extensive textual corpora. Although a direct strategy for preference learning involves supervised fine-tuning on exemplary human responses, the most effective methods to date belong to the category of reinforcement learning from human (or AI-generated) feedback (RLHF/RLAIF; \citep{christiano2017deep,bai2022constitutional}). RLHF techniques involve constructing a reward model trained on datasets of human choices and subsequently employing reinforcement learning to adjust the language model's outputs to yield high-reward responses, all while maintaining similarity to the pre-trained model. Despite RLHF yielding models that demonstrate strong abilities in conversation and code generation, its overall procedure is much more intricate compared to supervised approaches; it requires the training of several language models and frequent policy sampling throughout the training loop, which leads to substantial computational expenditures.

In this work, we demonstrate an approach for aligning language models with human preferences by optimizing directly against those preferences, bypassing the need for explicit reward modeling or the use of reinforcement learning techniques. We introduce \textit{{Direct Preference Optimization} (DPO)}, an algorithm designed to implicitly maximize the same reward (under a KL-divergence constraint) targeted by current RLHF strategies, while offering a more straightforward implementation and training process. In essence, DPO updates the model to favor responses that align with human preferences by raising the log-probability ratio between favored and disfavored outputs. Notably, it incorporates an adaptive, example-specific importance weighting that mitigates the model collapse often encountered with simpler probability ratio objectives. DPO shares with prior approaches the reliance on theoretical preference frameworks—such as the Bradley-Terry model \cite{bradley1952rankanalysis}—for quantifying agreement between reward functions and observed human choices. Unlike standard pipelines, which leverage this preference modeling to supervise reward function learning before optimizing policies against these learned rewards, DPO utilizes a change of variables to express the preference loss directly in terms of the policy itself. With a dataset composed of pairwise human preference judgments over model-generated responses, DPO enables policy optimization via a straightforward binary cross-entropy loss, effectively yielding a policy that is optimal with respect to an implicit reward derived from the collected preference information.

The core contribution of this work is Direct Preference Optimization (DPO), an uncomplicated algorithm that enables the training of language models from preference data without relying on reinforcement learning. Experimental results indicate that DPO performs on par with, or surpasses, current approaches—including PPO-based RLHF—when applied to preference-based tasks like sentiment control, text summarization, and conversational modeling, utilizing language models containing as many as 6B parameters.

\section{Related Work}

Progress in self-supervised language models of larger scales has enabled them to address certain tasks without any prior examples (zero-shot) \citep{radford2019language} or by conditioning on just a handful of demonstrations (few-shot) \citep{gpt3,megatron,chowdhery2022palm}. Nevertheless, performance on practical applications and alignment with user intent are substantially enhanced by further fine-tuning on instruction datasets paired with human-written responses \citep{mishra-etal-2022-cross,sanh2022multitask,chung2022scaling,thoppilan2022lamda}. This process, termed 'instruction-tuning,' equips LLMs with the ability to generalize to new instructions beyond those seen during tuning and generally boosts their practicality and flexibility \citep{chung2022scaling}. While instruction tuning has proven highly effective, it is frequently simpler to amass \textit{comparative} human judgments about output quality than it is to curate high-quality expert demonstrations. Consequently, later research has focused on further adapting LLMs using datasets based on human preference rankings, leading to improvements in areas such as translation \citep{kreutzer-etal-2018-reliability}, summarization \citep{stiennon2022learning,ziegler2020finetuning}, narrative generation \citep{ziegler2020finetuning}, and instruction following \citep{ouyang2022training,ramamurthy2023is}. The common approach involves first fitting a neural reward model to match the collected preference data, often utilizing a framework such as the Bradley-Terry model \citep{bradley1952rankanalysis}, and then fine-tuning the LLM using reinforcement learning methods—typically REINFORCE \citep{williams1992reinforce}, proximal policy optimization (PPO; \cite{schulman2017proximal}), or their extensions \citep{ramamurthy2023is}—to optimize for the learned reward. Related research employs instruction-tuned LLMs with human feedback to generate synthetic preference datasets for qualities like safety or harmlessness \citep{bai2022constitutional}, relying on limited human supervision delivered through textual rubrics that guide LLM-generated annotations. This set of methodologies integrates advances from two research streams: one dedicated to using reinforcement learning in training language models for various tasks~\citep{Ranzato2015SequenceLT,paulus2018a,wu2018learning}, and another focused on general frameworks for learning from human feedback \citep{christiano2017deep,kupcsik2018learning}. Despite the advantages of using preference-based feedback, reinforcement learning for fine-tuning large language models introduces substantial practical complexities; the present work offers a theoretically principled alternative for optimizing with relative preferences that sidesteps the need for RL.

Beyond language-focused applications, the problem of learning policies from preferences has been explored within both bandit and reinforcement learning paradigms, yielding a variety of methodologies. One such method, contextual dueling bandits (CDB; \cite{yue2012karmed,dudik2015contextual}), frames contextual bandit learning in terms of action preferences or rankings instead of explicit reward signals. In settings without absolute rewards, the analysis of CDBs introduces the concept of a \textit{von Neumann winner}, describing a policy whose average probability of outperforming any alternative policy is no less than 50\% \citep{dudik2015contextual}. Notably, CDBs assume that preference feedback is provided in an online fashion, in contrast to the learning-from-human-preferences framework, which predominantly relies on a predetermined dataset of action pairs annotated with offline preference judgments \citep{yan2022human}. Preference-based RL (PbRL), likewise, operates using binary feedback from an \textit{unknown} 'scoring' mechanism in lieu of rewards \citep{BusaFekete2014,ruiz2023dueling}. Multiple strategies exist for PbRL, some of which leverage off-policy preference datasets, though they typically involve first inferring the underlying scoring or reward function before proceeding to policy optimization \citep{jain2013learning,BusaFekete2014,christiano2017deep,sadigh2017active,kupcsik2018learning}. Contrasting with these multi-step approaches, we introduce a single-stage policy learning technique that directly updates policies to accommodate the specified preferences.

\section{Preliminaries}\label{section:prelims}

We summarize the RLHF workflow as described in \citeauthor{ziegler2020finetuning}, with subsequent refinements in works such as \citep{stiennon2022learning, bai2022training, ouyang2022training}. This framework is generally structured into three key steps: (1) supervised fine-tuning (SFT); (2) gathering preferences and training a reward model; and (3) optimizing using reinforcement learning.

\textbf{SFT}: The RLHF approach typically starts with the supervised fine-tuning of a pre-trained language model, leveraging carefully curated datasets corresponding to target tasks (such as dialogue systems or summarization). This process produces an initial model denoted $\pi^\text{SFT}$.

\textbf{Reward Modelling Phase}: During the second stage, the SFT model is tasked with generating answer pairs $(y_1, y_2)\sim \pi^\text{SFT}(y \mid x)$ in response to input prompts $x$. These answer pairs are subsequently evaluated by human annotators, who indicate a preference between the two, formalized as $y_w\succ y_l \mid x$, with $y_w$ representing the favored and $y_l$ the less favored option within $(y_1, y_2)$. It is postulated that such preferences arise from an underlying, unobserved reward function $r^*(y, x)$. To model these preference data, several methodologies can be adopted, with the Bradley-Terry (BT) model \cite{bradley1952rankanalysis} being widely utilized (though more sophisticated models, such as the Plackett-Luce ranking models \citep{plackett1975analysis, luce2012individual}, may be incorporated if rankings over more than two responses are available). Within the BT framework, the distribution over human preferences $p^*$ can be expressed as:

\begin{equation}\label{eq:bradley-terry}
    p^*(y_1\succ y_2 \mid x)=\frac{\exp\left(r^*(x, y_1)\right)}{\exp\left(r^*(x, y_1)\right) + \exp\left(r^*(x, y_2)\right)}.
\end{equation}

Given a fixed collection of comparisons $\mathcal{D}=\bigl\{x^{(i)}, y_w^{(i)}, y_l^{(i)}\bigr\}_{i=1}^N$ drawn according to $p^*$, we can introduce a reward model $r_{\phi}(x, y)$ and learn its parameters by employing maximum likelihood estimation. By recasting this task as a binary classification problem, the corresponding negative log-likelihood loss function becomes:

\begin{equation}\label{eq:reward_model}
    \mathcal{L}_R(r_{\phi}, \mathcal{D}) = -\mathbb{E}_{(x, y_w, y_l)\sim \mathcal{D}}\bigl[\log \sigma(r_{\phi}(x, y_w)- r_{\phi}(x, y_l))\bigr]
\end{equation}

where $\sigma$ denotes the logistic function. For LMs, the reward model $r_{\phi}(x, y)$ is typically initialized using parameters from the SFT model $\pi^\text{SFT}(y \mid x)$, augmented with an additional linear layer atop the final transformer layer to output a scalar value representing the reward \cite{ziegler2020finetuning}. To stabilize the reward predictions, previous studies apply normalization to the reward outputs such that $\mathbb{E}_{x,y\sim \mathcal{D}}\left[r_\phi(x, y)\right] = 0$ holds across all $x$.

\textbf{RL Fine-Tuning Phase}: In this stage, the language model receives feedback using the learned reward function. Consistent with methodologies in earlier works~\citep{jaques2017sequence, jaques2020human}, the training objective is posed as

\begin{equation}\label{eq:RL}
\max_{\pi_{\theta}}  \mathbb{E}_{x\sim \mathcal{D}, y\sim \pi_{\theta}(y \mid x)}\bigl[r_{\phi}(x, y)\bigr] - \beta\mathbb{D}_{\textrm{KL}}\bigl[\pi_{\theta}(y\mid x)\mid \mid \pi_\text{ref}(y\mid x)\bigr],
\end{equation}

Here, $\beta$ serves as a hyperparameter that regulates how much the learned policy can diverge from a base reference policy $\pi_\text{ref}$, specifically the original SFT model $\pi^\text{SFT}$. In real-world implementations, the policy $\pi_\theta$ used in the language model is also set to start from $\pi^\text{SFT}$. Incorporating this constraint is essential to keep the policy from straying excessively from the data distribution where the reward model operates reliably, in addition to preserving the diversity of generated outputs and averting the collapse of responses into a narrow set of high-reward options. Given that language generation operates over discrete outputs, the resulting objective is not differentiable and, therefore, is most often addressed with reinforcement learning techniques. The conventional method \citep{ziegler2020finetuning, stiennon2022learning, bai2022training, ouyang2022training} defines the reward function as ${r(x, y) = r_{\phi}(x, y) -\beta (\log \pi_{\theta}(y\mid x) - \log \pi_\text{ref}(y\mid x))}$, with maximization typically performed through PPO \cite{schulman2017proximal}.

\section{Direct Preference Optimization}\label{sec:DPO}

Prompted by the difficulties inherent in deploying reinforcement learning algorithms to large-scale tasks like the fine-tuning of language models, our objective is to develop a straightforward method for policy optimization that operates directly from preference data. In contrast to conventional RLHF approaches, which first construct a reward function and subsequently perform reinforcement learning to maximize it, our method utilizes a specific parameterization of the reward model. This parameterization allows us to analytically extract the optimal policy in closed form, thereby eliminating the need for an RL training loop. As we elaborate further below, our central contribution lies in exploiting a theoretical correspondence between reward functions and their optimal policies. This correspondence permits the conversion of a loss over reward models into a loss defined over policies themselves. By employing this change-of-variables technique, we circumvent the necessity of learning an explicit, isolated reward function, all while maintaining compatibility with established models of human preferences, including the Bradley-Terry model. Ultimately, this framework combines both the language model and an implicit representation of the reward within a single policy network.

\textbf{Derivation of the DPO objective.} Our analysis begins by considering the RL objective defined in Eq.~\ref{eq:RL}, incorporating a general reward function $r$, as utilized in earlier studies. Building on the methods of previous research~\citep{peters2007reinforcement, peng2019advantage, korbak2022reinforcement, go2023aligning}, it is direct to demonstrate that the optimal policy for the KL-constrained reward maximization problem presented in Eq.~\ref{eq:RL} is expressed as:

\begin{equation}\label{eq:op_policy}
    \pi_r(y\mid x) = \frac{1}{Z(x)}\pi_\text{ref}(y\mid x)\exp\left(\frac{1}{\beta}r(x, y)\right),
\end{equation}

Here, $Z(x) =\sum_{y}\pi_\text{ref}(y\mid x)\exp\left(\frac{1}{\beta}r(x, y)\right)$ denotes the partition function. For a detailed derivation, refer to Appendix \ref{app:derivation1}. Despite employing the maximum likelihood estimate $r_{\phi}$ as an approximation of the true reward function $r^*$, computing the partition function $Z(x)$ remains computationally demanding \citep{korbak2022reinforcement, go2023aligning}, thereby limiting the practical use of this formulation. Nevertheless, by manipulating Eq.~\ref{eq:op_policy}, the reward function can be re-expressed in terms of its associated optimal policy $\pi_r$, the reference policy $\pi_\text{ref}$, and the unobserved partition function $Z(\cdot)$. To achieve this, we apply the logarithm to both sides of Eq.~\ref{eq:op_policy} and perform straightforward algebraic manipulations:

\begin{equation}\label{eq:main_eq}
    r(x,y) =\beta \log \frac{\pi_r(y\mid x)}{\pi_\text{ref}(y\mid x)} + \beta \log Z(x).
\end{equation}

This reparameterization framework can be utilized for both the true reward function $r^*$ and its associated optimal policy $\pi^*$. Notably, in the Bradley-Terry preference model, the outcome depends exclusively on the reward difference between two completions: ${p^*(y_1 \succ y_2 \mid x) = \sigma(r^*(x, y_1) - r^*(x, y_2))}$. When substituting the expression for $r^*(x,y)$ from Eq.~\ref{eq:main_eq} into the Bradley-Terry formulation in Eq.~\ref{eq:bradley-terry}, the partition function terms eliminate each other. As a result, the probability of human preference can be articulated solely using the optimal policy $\pi^*$ and the reference policy $\pi_\text{ref}$. Therefore, under the Bradley-Terry framework, the optimal RLHF policy $\pi^*$ is characterized by the following preference structure:

\begin{equation}\label{eq:objective}
    p^*(y_1\succ y_2 \mid x)=\frac{1}{1 + \exp\left(\beta \log \frac{\pi^*(y_2\mid x)}{\pi_\text{ref}(y_2\mid x)} - \beta \log \frac{\pi^*(y_1\mid x)}{\pi_\text{ref}(y_1\mid x)}\right)}
\end{equation}

The full derivation can be found in Appendix~\ref{app:derivation2}. Although Eq.~\ref{eq:objective} relies on the Bradley-Terry model, analogous results can be developed for the broader family of Plackett-Luce models~\citep{plackett1975analysis, luce2012individual}; see Appendix~\ref{app:plackett_luce_models} for details.

Having expressed the likelihood of human preference data in terms of the optimal policy instead of the reward function, we can define a maximum likelihood objective with respect to a parameterized policy $\pi_\theta$. Mirroring the structure of the reward modeling formulation (as in Eq.~\ref{eq:reward_model}), this gives rise to the following policy objective:

\begin{equation}\label{eq:optimum_model}
    \mathcal{L}_\text{DPO}(\pi_{\theta}; \pi_\text{ref}) = -\mathbb{E}_{(x, y_w, y_l)\sim \mathcal{D}}\left[\log \sigma \left(\beta \log \frac{\pi_{\theta}(y_w\mid x)}{\pi_\text{ref}(y_w\mid x)} - \beta \log \frac{\pi_{\theta}(y_l\mid x)}{\pi_\text{ref}(y_l\mid x)}\right)\right].
\end{equation}

In this approach, we estimate an implicit reward through an alternative parameterization, where the corresponding optimal policy reduces to $\pi_\theta$. Additionally, our methodology is mathematically analogous to learning a reparameterized Bradley-Terry model, thereby inheriting various theoretical guarantees such as consistency, assuming appropriate conditions on the underlying preference data distribution \cite{bong2022generalized}. Section~\ref{sec:theory} further elaborates on DPO's theoretical aspects and its relationship to existing literature.

\textbf{How does the DPO update function?} To better understand the workings of DPO at a mechanistic level, we examine the gradient of its loss function, $\mathcal{L}_\text{DPO}$. Specifically, the derivative with respect to the parameter vector $\theta$ is given by:

\begin{multline*}\label{eq:gradient}
    \nabla_\theta \mathcal{L}_\text{DPO}(\pi_\theta;\pi_\text{ref}) = \\ -\beta\mathbb{E}_{(x, y_w, y_l) \sim \mathcal{D}} \bigg[\underbrace{\sigma(\hat{r}_\theta(x, y_l) - \hat{r}_\theta (x, y_w))}_\text{higher weight when reward estimate is wrong}\bigg[\underbrace{\nabla_\theta\log \pi(y_w \mid x)}_\text{increase likelihood of $y_w$} - \underbrace{\nabla_\theta\log\pi(y_l \mid x)}_\text{decrease likelihood of $y_l$}\bigg]\bigg],
\end{multline*}

where $\hat{r}_\theta(x, y) = \beta \log \frac{\pi_\theta(y \mid x)}{\pi_\text{ref}(y \mid x)}$ represents the reward function implicitly specified by the language model $\pi_\theta$ and a reference model $\pi_\text{ref}$ (see Section~\ref{sec:theory} for further details). In essence, the loss gradient for $\mathcal{L}_\text{DPO}$ serves to promote the probability of generating the preferred completions $y_w$, while suppressing the probability assigned to less preferred completions $y_l$. Notably, each example is weighted according to the extent that the implicit reward function $\hat{r}_\theta$ (scaled by $\beta$) overvalues the dispreferred completions; this quantifies the degree to which the ranking of completions by the implicit reward diverges from the target order and reflects the influence of the KL regularization parameter. Our empirical findings indicate that such weighting is crucial—removing the weighting factor from this approach, as in a simplistic baseline, often results in model collapse (see Appendix Table~\ref{tab:unlikelihood_generations}).

\textbf{DPO outline.} The typical workflow for DPO involves: 1) Generating candidate responses $y_1, y_2 \sim \pi_\text{ref}(\cdot \mid x)$ for each input $x$, then annotating these with human preference judgments to form an offline dataset $\mathcal{D} = \{x^{(i)}, y_w^{(i)}, y_l^{(i)}\}_{i=1}^N$; and 2) updating the language model $\pi_\theta$ by minimizing $\mathcal{L}_\text{DPO}$ with respect to the specified $\pi_\text{ref}$, dataset $\mathcal{D}$, and chosen $\beta$. In practical applications, it is often desirable to leverage existing public datasets containing human preferences, instead of synthesizing completions and collecting new human annotations. As these datasets are created with $\pi^\text{SFT}$ as the generative model, we set $\pi_\text{ref} = \pi^\text{SFT}$ where such a model is accessible. If $\pi^\text{SFT}$ is not accessible, we estimate $\pi_\text{ref}$ by maximizing the likelihood of preferred completions ${(x, y_w)}$, specifically by solving ${\pi_\text{ref} = \argmax_{\pi}\mathbb{E}_{x, y_w \sim \mathcal{D}}\left[\log \pi(y_w \mid x)\right]}$. Adopting this approach lessens the impact of distributional discrepancies between the unavailable true reference distribution and the $\pi_\text{ref}$ adopted within DPO. For more detailed information regarding the technical implementation and choice of hyperparameters, see Appendix~\ref{app:implementation}.

\section{Theoretical Analysis of DPO} This section offers a deeper theoretical interpretation of the DPO approach, substantiates it with formal justification, and discusses how DPO addresses certain challenges present in RLHF actor-critic frameworks, such as those encountered with PPO~\cite{schulman2017proximal}.

\label{sec:theory}

\subsection{Your Language Model Is Secretly a Reward Model} DPO manages to circumvent both the need for explicit reward fitting and separate RL-based policy learning, by employing a unified maximum likelihood objective. The optimization criterion in Eq. \ref{eq:main_eq} can actually be seen as a Bradley-Terry model, where the reward is parameterized by $r^*(x, y) = \beta \log\frac{\pi^*_\theta(y \mid x)}{\pi_\text{ref}(y \mid x)}$. Our training process focuses on directly optimizing the parameters of $\pi_{\theta}$, which corresponds to reward model training as formulated in Eq. \ref{eq:reward_model}, once a variable substitution is made. This section develops the theoretical underpinnings for this form of reparameterization, demonstrating that it does not restrict the space of reward models that can be learned, and that it admits perfect recovery of the optimal policy. We start by introducing a formal equivalence relation among reward functions.

\begin{definition}
Two reward functions $r(x, y)$ and $r'(x, y)$ are said to be equivalent if and only if there exists a function $f$ such that ${r(x, y) - r'(x, y) = f(x)}$.
\end{definition}

This relation is straightforwardly verified to be an equivalence relation, thus dividing the collection of reward functions into distinct equivalence classes. The following two lemmas can then be established:

\begin{lemma}\label{lemma:same_prefrence}
    Within the Plackett-Luce preference model, including the special case of Bradley-Terry, any pair of reward functions that belong to the same equivalence class generate identical preference distributions.
\end{lemma}

\begin{lemma}\label{lemma:same_policy}
    For the constrained RL problem, if two reward functions are members of the same equivalence class, they will result in identical optimal policies.
\end{lemma}

The demonstrations are elementary and can be found in Appendix \ref{app:lemma1}. The initial lemma highlights a commonly recognized under-specification problem associated with the Plackett-Luce model family \cite{plackett1975analysis}. Owing to this ambiguity, it becomes necessary to impose supplementary identifiability constraints in order to secure meaningful guarantees for the MLE estimates as formulated in Eq. \ref{eq:reward_model} \cite{bong2022generalized}. The subsequent lemma establishes that all reward functions within the same equivalence class induce identical optimal policies; consequently, our ultimate aim is to recover any representative reward function from this optimal class. The subsequent Theorem is formally established in Appendix~\ref{app:thm1}:

\begin{theorem}\label{thm:main}
    Assuming certain mild conditions, every reward class that aligns with the Plackett-Luce (and specifically the Bradley-Terry) models can be expressed using the following reparameterization: ${r(x, y) = \beta \log \frac{\pi(y\mid x)}{\pi_\text{ref}(y\mid x)}}$, where $\pi(y\mid x)$ denotes some model, and $\pi_\text{ref}(y \mid x)$ is a fixed reference model.
\end{theorem}

\begin{sproof}
    Take an arbitrary reward function $r(x, y)$, which yields the optimal model $\pi_r(y \mid x)$ as defined in Eq. \ref{eq:op_policy}. We aim to demonstrate that any reward function belonging to the same equivalence class as $r$ can be expressed using the previously introduced reparameterization. We construct the projection $f$ as follows:
\begin{equation}
    f(r; \pi_\text{ref}, \beta)(x, y) = r(x, y) - \beta\log\sum_{y}\pi_\text{ref}(y\mid x)\exp\left(\frac{1}{\beta}r(x, y)\right)
\end{equation}
Here, the operator $f$ effectively standardizes the reward function by subtracting the logarithm of $\pi_r$'s partition function, scaled by $\beta$. Since the normalization term depends solely on $x$ (i.e., the prefix), $f(r; \pi_\text{ref}, \beta)(x, y)$ constitutes another reward function in the same equivalence class as $r(x, y)$. Substituting the right-hand side of Eq.~\ref{eq:main_eq} for $r$ (valid for any reward function), we find that $f(r; \pi_\text{ref}, \beta)(x, y) = \beta \log \frac{\pi_r(y\mid x)}{\pi_\text{ref}(y\mid x)}$. Therefore, the projection $f$ generates an element of $r$'s equivalence class in the specified form, showing that the proposed reparameterization does not restrict the expressiveness of our reward model.
\end{sproof}

Theorem~\ref{thm:main} can also be interpreted as providing an explicit selection criterion for the reward function from each equivalence class in the context of the DPO reparameterization, namely, the reward function that fulfills:

\begin{equation}\label{eq:lag_p}
     \sum_{y}\underbrace{\pi_\text{ref}(y\mid x)\exp\left(\frac{1}{\beta}r(x, y)\right)}_{=\pi(y\mid x)\text{, using Thm.~\ref{thm:main} reparam.}} = 1,
\end{equation}

That is, $\pi(y\mid x)$ forms a legitimate probability distribution, with all entries positive and summing to unity. By referring to Eq.~\ref{eq:op_policy}, we recognize Eq.~\ref{eq:lag_p} as representing the partition function associated with the optimal policy that the reward function $r(x, y)$ defines. The central contribution of the DPO method is the imposition of well-chosen constraints upon the otherwise underdetermined Plackett-Luce (and particularly Bradley-Terry) preference model families. Through these constraints, the space of feasible reward models remains unchanged, but the induced optimal policy in Eq.~\ref{eq:op_policy} becomes analytically manageable for any prompt $x$.

\subsection{Instability of Actor-Critic Algorithms} Our framework can further be leveraged to analyze the sources of instability inherent in commonly employed actor-critic methods for RLHF, including PPO. We adopt the typical RLHF setup and center our discussion on the RL fine-tuning phase described in Section \ref{section:prelims}. This approach allows us to relate the underlying mechanics to the control as inference paradigm \cite{levine2018reinforcement}, specifically in the context of the constrained RL formulation in \ref{eq:RL}. Given a parameterized policy $\pi_{\theta}(y\mid x)$, we aim to minimize the divergence $\mathbb{D}_{\text{KL}}[\pi_{\theta}(y|x) \mid \mid \pi^*(y\mid x)]$, where $\pi^*$, defined in Eq. \ref{eq:optimum_model}, represents the optimal policy derived from the reward model $r_{\phi}(y, x)$. Through some mathematical manipulation, this leads us to the following optimization criterion:

\begin{equation}\label{eq:AC}
    \max_{\pi_{\theta}}\mathbb{E}_{\pi_{\theta}(y\mid x)}\bigg[\underbrace{r_{\phi}(x, y) -\beta\log\sum_{y}\pi_\text{ref}(y\mid x)\exp\left(\frac{1}{\beta}r_{\phi}(x, y)\right)}_{f(r_{\phi}, \pi_\text{ref}, \beta)} - \underbrace{\beta\log\frac{\pi_{\theta}(y\mid x)}{\pi_\text{ref}(y\mid x)}}_{\text{KL}}\bigg]
\end{equation}

This objective is identical to that optimized in previous studies 
\citep{ziegler2020finetuning, stiennon2022learning, bai2022training, ouyang2022training}, where the reward associated with $r_{\phi}$ follows the DPO-equivalent formulation. Within this context, the normalization component of $f(r_{\phi}, \pi_\text{ref}, \beta)$ can be seen as the soft value function corresponding to the reference policy $\pi_\text{ref}$. Although this normalization does not alter the location of the optimum, omitting it can lead to policy gradient estimates with increased variance, which, in turn, may destabilize training. This normalization effect can be incorporated through the use of a learned value function, but such an approach can present optimization challenges. Prior approaches have also addressed normalization by using a human completion as a baseline, effectively approximating the normalization term with a single-sample Monte-Carlo estimate. In contrast, the DPO reparameterization results in a reward function that eliminates the necessity for such baselines.

\section{Experiments}
Here, we present an empirical assessment of DPO’s effectiveness in learning policies from preference data. We begin by investigating DPO in a controlled text-generation environment, focusing on the efficiency of its balance between reward optimization and KL-divergence regularization with the reference policy. This comparison is made against established preference-based training approaches, including PPO. Subsequently, we assess DPO’s capabilities on more complex RLHF problems, such as dialogue and summarization, using larger-scale models. Our findings indicate that, with minimal hyperparameter tuning, DPO generally matches or surpasses robust baselines like RLHF with PPO, and also outperforms selecting the best among $N$ sampled trajectories scored by a learned reward model. Prior to discussing the outcomes, we outline our experimental methodology; supplementary information is available in Appendix~\ref{app:exp_details}.

\textbf{Tasks.} We conduct experiments on three distinct open-ended text generation tasks. Across all tasks, learning algorithms derive a policy from a dataset of preference tuples $\mathcal{D}=\bigl\{x^{(i)}, y_w^{(i)}, y_l^{(i)}\bigr\}_{i=1}^N$. For \textbf{controlled sentiment generation}, each $x$ consists of a movie review prefix sourced from the IMDb dataset \cite{maas-EtAl:2011:ACL-HLT2011}, with the objective for the policy to generate a continuation $y$ that expresses positive sentiment. To ensure rigorous evaluation, we synthetically create preference pairs based on generations, leveraging a pre-trained sentiment classifier such that $p(\text{positive}\mid x,y_w)>p(\text{positive}\mid x,y_l)$. The SFT approach involves fine-tuning GPT-2-large on the training portion of the IMDB reviews until it converges (see App~\ref{app:sentiment_details} for more on methodology). In the \textbf{summarization} task, $x$ represents a Reddit forum post, and the policy aims to produce a concise summary $y$ capturing the main ideas. We adopt the Reddit TL;DR summarization dataset \citep{volske-etal-2017-tl}, complemented by human preference data curated by \citeauthor{stiennon2022learning}. For supervised fine-tuning (SFT), we utilize a model adapted to human-created Reddit summaries\footnote{\url{https://huggingface.co/CarperAI/openai_summarize_tldr_sft}}, implemented via the TRLX \citep{leandro_von_werra_2023_7790115} toolkit for RLHF. Notably, the preferences used are annotated on outputs from another similarly fine-tuned SFT model, as detailed by \citeauthor{stiennon2022learning}. Lastly, for the \textbf{single-turn dialogue} task, $x$ is a user prompt spanning diverse domains, such as physics questions or requests for personal advice. The policy must return a relevant and useful response $y$. Our study uses the Anthropic Helpful and Harmless dialogue dataset \citep{bai2022training}, which includes 170k conversational exchanges between people and an automated assistant. Each dialogue transcript is finalized with two alternative model-generated replies and a human judgment indicating the preferred one. As there is no publicly available SFT model for this data, we construct the SFT model by fine-tuning a publicly available language model solely on preferred response samples.

\textbf{Evaluation.} Our experimental design incorporates two main evaluation strategies. To assess each algorithm’s capability in constrained reward maximization, we perform a controlled sentiment generation task. Here, we plot the trade-off frontier between the attained reward and the KL-divergence relative to the reference policy; because the actual reward function (given by a sentiment classifier) is available in this setup, this analysis is feasible. Conversely, practical deployments lack access to the true reward function. Thus, we rely on \textit{win rate} comparisons against a baseline policy, using GPT-4 to approximate human judgment of quality in both summarization and single-turn dialogue contexts. In summarization tasks, we select reference summaries from the test set as baselines, while for dialogue we utilize the human-preferred response in the test data. Prior research indicates that language models may outperform conventional metrics as automatic evaluators \citep{Chen2023ExploringTU}. Nevertheless, to validate the reliability of GPT-4 in our evaluation, we additionally conduct a human user study, detailed in Sec.~\ref{sec:human-judgments}. Our findings reveal that GPT-4’s evaluations align closely with human assessments, with agreement levels between humans and GPT-4 matching or exceeding typical inter-annotator reliability.

\textbf{Methods.} Beyond DPO, we assess a range of established techniques designed to align language model outputs with human preferences. For the summarization task, we utilize zero-shot prompting using \textbf{GPT-J} \citep{gpt-j}, and for the dialogue task, we apply 2-shot prompting with \textbf{Pythia-2.8B} \citep{biderman2023pythia}. Additionally, we examine the performance of the \textbf{SFT} model and \textbf{Preferred-FT}—the latter being a variant obtained by supervised fine-tuning on the selected completion $y_w$ produced by either the SFT model (in the controlled sentiment and summarization scenarios) or by a generic LM (for the single-turn dialogue task). Another related approach is \textbf{Unlikelihood}~\citep{welleck2019neural}, which trains the policy not only to increase the likelihood of $y_w$ but also to actively reduce the likelihood of $y_l$; an adjustable weighting parameter $\alpha\in[0,1]$ modulates the influence of the ‘unlikelihood’ component. We further include \textbf{PPO} \citep{schulman2017proximal}, which leverages a reward model learned from preference data, and its oracle variant, \textbf{PPO-GT}, which has access to the true reward function under the controlled sentiment condition. Our sentiment-related studies employ two variants of PPO-GT: a standard implementation \cite{leandro_von_werra_2023_7790115} and a customized version that incorporates reward normalization and additional hyperparameter tuning for enhanced results (these refinements are also used when applying PPO with learned reward signals). Lastly, we evaluate the \textbf{Best of $N$} strategy, which involves generating $N$ candidate outputs from the SFT (or Preferred-FT for dialogue), and selecting the one that achieves the highest score according to a reward model trained on preference data. While this approach can yield high-quality responses by isolating the reward model's influence from PPO training, it is generally infeasible for moderate $N$ due to the requirement of generating $N$ samples per input at inference time.

\subsection{How well can DPO optimize the RLHF objective?}

The reward maximization framework with a KL constraint, commonly employed in RLHF algorithms, seeks to maximize reward while simultaneously limiting divergence from the reference policy. Consequently, when evaluating and comparing algorithms, it is essential to consider both the attained reward and the associated KL divergence; obtaining marginally higher rewards at the expense of significantly greater KL divergence may not be optimal. Figure~\ref{fig:frontier-tldr-main} depicts the reward-KL tradeoff frontier for several algorithms within the sentiment domain. For each algorithm, we conduct several training runs, each with a distinct setting for policy conservativeness (target KL values of $\{3,6,9,12\}$ for PPO, $\beta \in \{0.05,0.1,1,5\}$, $\alpha\in\{0.05,0.1,0.5,1\}$ for unlikelihood, and varying random seeds for preferred-FT). In total, this results in 22 different experimental runs. Following every 100 training updates up to convergence, the resulting policies are evaluated on a suite of test prompts, measuring both the mean reward according to the true reward function and the average sequence-level KL\footnote{That is, the sum of the per-timestep KL-divergences.} relative to the reference policy $\text{KL}\left(\pi\mid \mid \pi_\text{ref}\right)$. Our experiments demonstrate that DPO yields a considerably more favorable frontier, securing the highest rewards while maintaining low KL divergence. This outcome is particularly striking for several reasons. Most notably, despite both DPO and PPO optimizing the same objective, DPO consistently exhibits greater efficiency; the reward/KL tradeoff achieved by DPO is strictly superior to that of PPO. Furthermore, DPO surpasses PPO even in cases where PPO is provided with access to the ground truth rewards (PPO-GT).

\subsection{Can DPO scale to real preference datasets?}
\label{sec:dpo-real-datasets}
We proceed to assess how effectively DPO can be fine-tuned on tasks involving summarization and single-turn dialogue, both of which rely on real human preference data. In the summarization domain, established automated evaluation criteria like ROUGE often fail to align well with human judgments~\citep{stiennon2022learning}, and earlier studies indicate that language models adjusted via PPO using human feedback can produce superior summaries. For our comparison, we generate completions using different methods on the test portion of the TL;DR summarization dataset, subsequently calculating the mean win rate when matched against the reference summaries in the same test set. To ensure a comprehensive evaluation, outputs are sampled across a range of temperatures from 0.0 to 1.0, and the respective win rates are visualized in Figure~\ref{fig:frontier-tldr-main} (right). All approaches—DPO, PPO, and Preferred-FT—begin with fine-tuning a shared GPT-J SFT model\footnote{\url{https://huggingface.co/CarperAI/openai_summarize_tldr_sft}}. Our analysis reveals that DPO achieves a win rate near 61\% at a sampling temperature of 0.0, outperforming PPO’s optimal result of roughly 57\% at the same temperature. Additionally, DPO attains a superior peak win rate compared to the best-of-$N$ baseline. It's important to mention that DPO's $\beta$ hyperparameter was not extensively optimized in our experiments, which suggests these findings could understate DPO’s full capabilities. Furthermore, DPO demonstrates greater resilience to variations in sampling temperature, whereas PPO's effectiveness can drop significantly at higher temperatures, approaching the level of the initial GPT-J model. Preferred-FT yields minimal improvement over the SFT baseline. Lastly, as detailed in Section~\ref{sec:human-judgments}, we conduct direct human assessments of DPO and PPO, where DPO outputs sampled at temperature 0.25 are selected over PPO outputs at temperature 0 by human annotators in 58\% of cases.

For single-turn dialogue evaluation, we assess the various approaches using the subset of the Anthropic HH dataset \citep{bai2022training} test split that features a single round of human-assistant interaction. In our GPT-4-based evaluations, we use the preferred completions from the test set as the reference point, calculating win rates across different techniques. Given the absence of a standard SFT baseline for this benchmark, our process begins with a pre-trained Pythia-2.8B model; we then apply Preferred-FT to adapt the reference model to selected completions—ensuring these outputs remain within the model's distribution—before proceeding to train using DPO. We also benchmark against the strongest completion out of 128 Preferred-FT samples (noting, as shown in Appendix Figure~\ref{fig:best-of-n}, that performance saturates at 128 completions for this task) and include a 2-shot prompted configuration of the base Pythia-2.8B. Our analysis indicates that DPO either matches or surpasses other methods when using each technique’s optimal sampling temperatures. Additionally, we evaluate a PPO-based RLHF model trained on the Anthropic HH dataset \footnote{\url{https://huggingface.co/reciprocate/ppo_hh_pythia-6B}}, drawing on a reputable implementation \footnote{\url{https://github.com/CarperAI/trlx/tree/main/examples/hh}}, but fail to identify a prompt or temperature setting that achieves results exceeding those of the base Pythia-2.8B. Considering our TL;DR experiments and the observation that both PPO and Best of 128 maximize the same reward function, we take the latter as an approximate surrogate for PPO performance. In summary, DPO emerges as the sole computationally tractable method that outperforms the reference completions in the Anthropic HH dataset, and its results are comparable to or exceed those of the resource-intensive Best of 128 baseline. Moreover, as illustrated in Figure~\ref{fig:dialogue-main}, DPO attains its optimal results after relatively few training iterations.

\subsection{Generalization to a new input distribution}

To extend our analysis of PPO and DPO performance under distributional changes, we assess the policies trained in our Reddit TL;DR summarization study on an out-of-domain dataset: the test split of CNN/DailyMail news articles \citep{nallapati-etal-2016-abstractive}. We employ the optimal sampling temperatures found for TL;DR (0 and 0.25) in these evaluations. Table~\ref{tab:ood} displays the resulting metrics. For evaluation, we measured the GPT-4 win rate relative to ground-truth summaries, utilizing the identical GPT-4 (C) prompt as in the Reddit TL;DR setting, substituting ``forum post'' with ``news article''. On this out-of-distribution task, DPO maintains a clear advantage over PPO. These findings offer preliminary support that DPO policies may achieve a degree of generalization on par with PPO policies, despite the fact that DPO does not rely on the extra unlabeled Reddit TL;DR prompts leveraged by PPO.

\subsection{Validating GPT-4 judgments with human judgments}
\label{sec:human-judgments}
To assess the trustworthiness of GPT-4's evaluation capabilities, we run a human annotation study leveraging outcomes from the TL;DR summarization task alongside two GPT-4 prompt variants. The first, \textbf{GPT-4 (S)} (simple), asks the model to identify which summary more effectively captures the essential content of a post. The second, \textbf{GPT-4 (C)} (concise), extends this by also querying which summary is more concise; this variation is included because, when using the \textbf{GPT-4 (S)} prompt, we observe a bias in GPT-4 favoring lengthier and more redundant summaries compared to human raters. Full prompt texts are provided in Appendix~\ref{app:prompts}. We structure our analysis around three pairwise comparisons, choosing systems with highest (DPO, temp. 0.25), lowest (PPO, temp. 1.0), and intermediate (SFT, temp. 0.25) performance, in order to span a range of sample qualities; all methods are compared to PPO with greedy decoding at its optimal temperature. Our findings indicate that, for both prompts, GPT-4's decisions correlate with those of human raters to a degree comparable to inter-annotator agreement among humans, indicating that GPT-4 serves as a reliable stand-in for human judgment (as we have a limited pool of human raters, multiple independent human evaluations are gathered only for DPO and PPO-1 comparisons). On the whole, the \textbf{GPT-4 (C)} prompt yields win rates that align more closely with human preferences, leading us to select this prompt for the principal analyses in Section~\ref{sec:dpo-real-datasets}. Further information about the human evaluation procedure, such as the rating interface and participant roster, can be found in Appendix~\ref{app:human-study}.

\section{Discussion}
Preference-based learning presents a robust and scalable approach for developing capable and value-aligned language models. In this work, we have presented DPO, a straightforward methodology for training language models using preference data without relying on reinforcement learning techniques. Instead of adapting the preference learning objective to fit conventional RL frameworks and leveraging standard RL methods, DPO establishes a direct correspondence between language model policies and reward functions. This connection allows for aligning language models with human preferences through a basic cross-entropy loss function, bypassing the need for reinforcement learning and without any reduction in general applicability. Remarkably, DPO achieves comparable or superior results to current RLHF methods, including those employing PPO, with minimal adjustment of hyperparameters. Consequently, DPO substantially lowers the difficulty of training language models on human preferences.

\textbf{Limitations \& Future Work.} Our findings open up multiple avenues for subsequent investigation. One pressing question concerns how the DPO policy performs on out-of-distribution data relative to approaches that employ an explicit reward function. Preliminary evidence indicates that DPO policies exhibit generalization capabilities akin to those achieved with PPO-based systems, yet further, more exhaustive evaluations are warranted. For instance, it remains to be seen whether leveraging self-labeling from the DPO policy is as effective in utilizing unlabeled prompts. Another subject for exploration involves understanding the characteristics of reward over-optimization within the framework of direct preference optimization, as well as determining whether the modest reduction in performance observed in Figure~\ref{fig:dialogue-main}-right can be attributed to this phenomenon. Moreover, though our study examines models containing up to 6B parameters, extending DPO to much larger, cutting-edge architectures represents a promising research direction. With respect to assessment methods, our analysis reveals that GPT-4-derived win rates are sensitive to prompt phrasing; therefore, future investigations might focus on optimizing prompts to yield more reliable judgments from automated evaluators. Lastly, while DPO has been demonstrated in the context of language modeling aligned with human preferences, its application may extend to training generative models in additional domains.

\end{document}